\section{Related Work}\label{sec:related}

\subsection{Self-repair and conditional component importance}\label{sec:rw-conditional}

Transformer circuits can adapt when components are ablated. The IOI circuit analysis of \citet{wang2023ioi} identified primary and backup Name Mover heads in GPT-2 Small, providing an early concrete setting in which several attention heads participate in a behavior and some become especially relevant under intervention. Recent prompt-specific circuit work further shows that mechanisms can vary systematically across templates even within IOI, reinforcing the need to state the distribution over which a mechanistic claim is supported \citep{franco2026promptspecific}. Subsequent work established self-repair more broadly. The Hydra effect describes compensatory changes in downstream language-model computation after ablation \citep{mcgrath2023hydra}. Copy-suppression analysis showed that self-repair can arise through specific functional mechanisms rather than a single generic backup process \citep{mcdougall2023copysuppression}. \citet{rushing2024selfrepair} further found self-repair across model families and the full training distribution, with substantial prompt-level variability and contributions from mechanisms including LayerNorm scaling and anti-erasure.

FALI does not claim the discovery of self-repair or heterogeneous compensation. More directly, it also does not introduce the idea of conditional component importance. \citet{gong2026coax} formalize Conditional Co-Ablation (CoAx), which scores how much a remaining unit's ablation effect grows once a primary set has been removed. Their experiments show that this conditional score can recover backup heads missed by first-order rankings and transfers from GPT-2 IOI to induction settings across multiple models. This establishes both the empirical relevance and the practical utility of asking how component importance changes under intervention.

The distinction is therefore one of explanatory stage. Conditional ablation answers whether, or to what extent, a component becomes more important in a perturbed circuit. FALI assumes such a candidate has already been selected and asks which changes in the candidate's attention-score computation, value content, and surrounding context account for that cross-state change. In this sense CoAx and FALI are complementary: the former can identify \emph{which} components warrant conditional analysis, whereas FALI provides a structured mechanism-resolution experiment for a chosen attention head.

\subsection{Interaction-aware causal attribution}\label{sec:rw-interaction}

Activation patching and related interchange interventions are standard tools for testing whether internal activations causally contribute to model behavior. Path patching refines this idea to hypotheses about particular computational paths \citep{goldowskydill2023pathpatching}; systematic studies have also shown that patching conclusions depend on choices such as corruption, metric, and intervention design \citep{zhang2024activationpatching,heimersheim2024activationpatching}. These methods are primarily localization or causal-testing tools: they identify components, activations, or paths whose replacement changes an outcome. FALI uses the same broad experimental tradition but addresses a different estimand---the change in one selected head's output-grounded causal contribution between two circuit states.

A particularly close conceptual neighbor is \citet{vaidyanathan2026curse}. Re-deriving activation patching through causal mediation analysis, they show that a component's natural indirect effect can include interaction terms that measure how its effect depends on the state of other model components. In GPT-2 IOI these interactions can hide conditionally important components, inflate others, and contribute to instability in component rankings. This result rules out any novelty claim that FALI is the first method to observe context-dependent causal effects. Instead, it motivates the next question: once a conditional or interaction-dependent effect has been identified, can we resolve that change into architecture-level factors of the selected attention head?

Causal Head Gating provides a complementary head-level perspective. \citet{nam2025causalhead} learn gates over attention heads, classify their roles as facilitating, interfering, or irrelevant, and find that head roles can depend on interactions with other heads. FALI is neither a learned gating method nor a global head-role taxonomy. It fixes one candidate head and two model states, then intervenes on explicitly defined internal factors of that head and the surrounding model. Likewise, we do not claim that FALI is generally more faithful than patching: intervention-based interpretability results can be sensitive to methodology, as demonstrated by circuit-faithfulness studies \citep{miller2024faithfulness}. Our robustness experiments therefore support only the intervention variants actually tested.

These interaction-aware results also inform our interpretation of the context factor $C$. Under FALI, $C$ is not a localized ``downstream sensitivity'' variable. It denotes the surrounding model computation that differs between intact and perturbed contexts outside the explicitly injected candidate score/value state. It can therefore aggregate changes in other heads, MLPs, residual representations, normalization, and upstream or downstream computation. A large $\phi_C$ is evidence about the specified counterfactual game, not identification of a unique external mechanism.

\subsection{Decomposing attention computation}\label{sec:rw-attention}

The separation between attention routing and transported content is foundational rather than novel to FALI. The transformer-circuits framework distinguishes the query--key (QK) circuit, which determines attention scores and patterns, from the output/value (OV) circuit, which determines what information is written when a source position is attended to \citep{elhage2021mathematical}. Later causal work makes increasingly fine distinctions inside attention computation. \citet{franco2025attentioncausal} identify low-dimensional communication signals with a causal relationship to attention patterns, while \citet{wang2026fewshot} causally separate Query--Key routing from Value updates when explaining how few-shot examples shape function vectors. Recent factorial intervention work likewise shows that attention allocation and token content can be experimentally separated yet strongly coupled in a different robustness setting \citep{zhu2026lexical}.

FALI builds on these distinctions but changes the object being decomposed. Its dependent quantity is not the attention pattern itself, nor the quality of a function vector, but the \emph{cross-state change in the selected head's downstream causal usefulness}. SCORE and VALUE are therefore accompanied by an explicit CONTEXT factor. Within SCORE, FALI further partitions structurally valid score coordinates into task-target scores $T$, rival scores $R$ in the same task-relevant query row, and all other valid scores $O$. This target/rival/other partition is hypothesis-driven: it requires a task-relevant target edge or edge set to be specified in advance, and is not an automatic circuit-discovery procedure.

Because score, value, and context can interact non-additively, FALI uses the standard Owen value for a game with a priori coalition structure \citep{owen1977values}. The coalition structure is fixed as SCORE=$\{T,R,O\}$, VALUE=$\{V\}$, and CONTEXT=$\{C\}$. Owen attribution is established game-theoretic machinery; our use of it is motivated by the requirement that the fine score contributions reconcile with the coarse SCORE contribution rather than by a claim of a new attribution theorem. The resulting values should also be read relative to the intervention semantics. Hybrid internal states need not be states the unmodified model would naturally visit, and causal-intervention research has documented that representation-level interventions can move models away from their natural internal distributions \citep{grant2025divergent}. We therefore treat FALI as a controlled counterfactual decomposition, not as evidence that the model implements five independent natural pathways.

Finally, our emphasis on explicit validation reflects a broader evaluation problem in mechanistic interpretability. Benchmarks such as MIB prioritize precise recovery of causal pathways or variables rather than plausible narrative alone \citep{mueller2025mib}. FALI's synthetic single-factor systems and implementation invariants serve the narrower purpose of checking that the intervention pipeline recovers known manipulated factors before empirical use. They do not remove the substantive limitations of target selection, coalition choice, or intervention dependence, which we address explicitly in Section~\ref{sec:limitations}.

An earlier public preprint introduced the name FALI for \emph{Frozen Attention Logit Intervention} and analyzed target-versus-background attention-logit changes in GPT-2 IOI \citep{sharmila2026attentionrepairs}. The present work supersedes that attention-probability interpretation. It changes the estimand to output-grounded conditional causal usefulness; expands the intervention game to $T/R/O/V/C$; uses standard Owen attribution with coarse--fine reconciliation; adds controlled implementation validation and broad canonical confirmation; and includes a prospectively frozen cross-setting transfer test whose fine-mechanism hypothesis fails. We therefore treat the earlier method as a nested attention-probability special case rather than as evidence for the present causal mechanism claims.
